% Minimal EigenEngine manuscript skeleton: REVTeX4-2 base + the eigen house layer.
% Compile: pdflatex main.tex  (or: latexmk -pdf main.tex)
\documentclass[reprint,aps,prx,superscriptaddress,nofootinbib]{revtex4-2}
\usepackage{eigen-revtex}

\begin{document}

\title{An EigenEngine Manuscript Skeleton}
\author{The Laplace Engine (Eigen)}
\affiliation{The Laplace Engine}
\author{G.\ McCaul}
\affiliation{The Laplace Engine}

\begin{abstract}
This skeleton exercises the eigen house layer on the REVTeX4-2 base: the shared
palette, the phys-coloured links, and the \verb|\claim| provenance tag. It is the
net-new starting point for a projection, and the fixed target for the template gate.
\end{abstract}

\maketitle

\section{Ground}\label{sec:ground}
The engine states a claim and resolves it to an artifact.\claim{1} Section~\ref{sec:ground}
is its own referent, so the reference is defined and the log stays clean. Inline maths
such as $\hat{H}\ket{\psi} = E\ket{\psi}$ uses the canonical notation stack.

\section{Outlook}\label{sec:outlook}
The closing section follows the author's convention and is named Outlook, never
Conclusion.

\begin{thebibliography}{9}
\bibitem{ref1} A.\ Author, \emph{A Title}, J.\ Ref.\ \textbf{1}, 1 (2026).
\end{thebibliography}

\end{document}
